\documentclass{article}
\usepackage[T1]{fontenc}
\usepackage{lmodern}
\usepackage{graphicx}
\usepackage{amsmath,amssymb}
\usepackage[a4paper,margin=2cm]{geometry}
\usepackage[absolute,overlay]{textpos}
\setlength{\TPHorizModule}{1mm}
\setlength{\TPVertModule}{1mm}
\usepackage[indonesian]{babel}
\usepackage{hyperref}
\setlength{\emergencystretch}{2em}
% unit: Halaman 72

\begin{document}

% === Halaman 72 (llm) ===

Pesan yang akan dienkripsi diatur terlebih dahulu sebagai berikut (langkah \textit{preprocessing}):

\begin{enumerate}
  \item Buang semua spasi di dalam pesan
  \item Ganti huruf j (bila ada) dengan i. Contoh: jalan $\rightarrow$ ialan
  \item Tulis pesan dalam pasangan huruf (\textit{bigram}).
  \item Jangan sampai ada bigram dengan pasangan huruf yang sama. Jika ada, sisipkan x di tengahnya
  \item Jika jumlah huruf di dalam pesan ganjil, tambahkan huruf x pada bigram terakhir
\end{enumerate}

\end{document}
